\section{Proof Sketches and Technical Remarks}
\label{sec:proof-sketches}

This section provides concise technical remarks clarifying the logical status
and scope of the results established in
Section~\ref{sec:impossibility}.
The purpose is not to introduce additional assumptions, strengthen theorems, or
extend proofs, but to make explicit what follows strictly from the Ontology of
Continua (OC) and the frozen executable specification
\texttt{ETS.K\_EA.v1.1}, and what does not.

\subsection{Source of the Impossibility Results}

The fundamental structural fact underlying all impossibility statements is that
KPI-based observation is represented by a map
\[
\pi_K:\Omega_{\mathrm{obs}}\to\Sigma_{\mathrm{obs}},
\]
which necessarily factors through the quotient
\(\Omega_{\mathrm{obs}}/\!\sim_K\).

Once observation is expressed in this form, non-identifiability follows
immediately: any two distinct states lying in the same fiber of \(\pi_K\) are
indistinguishable by construction.
This argument is purely set-theoretic and logical.
No assumptions about smoothness, topology, dimensionality, or regularity of
\(\Omega_{\mathrm{obs}}\) are required.

In this sense, the impossibility results here are structurally analogous to
classical observation-induced ambiguities in control and identification theory,
where indistinguishability arises from projection onto a lower-information
output space (cf.\ \cite{kalman1960,ljung1999}).

\subsection{On Injectivity and Admissible Assumptions}

Many classical impossibility or non-identifiability results in system
identification are derived under auxiliary assumptions on the observation map,
such as genericity, excitation conditions, smoothness, or dimensional mismatch
between state and output spaces \cite{ljung1999,walter1997}.

In contrast, this preprint deliberately avoids importing any such assumptions as
new primitives.
Doing so would exceed the OC/ETS scope and undermine the diagnostic character of
the results.

Accordingly, all statements are formulated in the minimal logically correct
form:
\begin{itemize}[leftmargin=2em]
  \item if the KPI readout map \(\pi_K\) is not injective, then distinct
        KPI-indistinguishable states exist
        (Theorem~\ref{thm:nonident});
  \item OC and \texttt{ETS.K\_EA.v1.1} do not imply injectivity of any admissible
        finite KPI family, and therefore do not support claims of universal KPI
        completeness
        (Theorem~\ref{thm:no-complete});
  \item any deterministic transformation applied to KPI readouts preserves
        non-identifiability
        (Proposition~\ref{prop:no-pp}).
\end{itemize}

This framing reflects the intended scope: a formal diagnostic boundary, not an
engineering guideline or modeling prescription.

\subsection{Observable Domain and Loss of Observability}

The observable domain \(\Omega_{\mathrm{obs}}\) is a strict subset of the full
enterprise state space \(\Omega(K_{EA})\).
It is defined exclusively via ETS-level observability structures, namely the
existence of the factor \(S_{\mathrm{obs}}\) and the validity of the
\emph{ObservabilityInvariant}.

States for which the ETS phase condition \texttt{observability\_lost} applies
are explicitly outside the scope of this analysis.
In such states, KPI readouts are not assumed to be defined, and no claims about
identifiability or non-identifiability are made.
All impossibility results in this preprint are therefore conditional on
admissible observability.

This separation mirrors standard distinctions in control theory, where
questions of identifiability are only meaningful on domains where outputs are
well-defined \cite{kalman1960}.

\subsection{Why Sampling and Aggregation Cannot Help}

Any aggregation over time windows, smoothing, scoring, embedding, or similar
analytic operation applied to KPI values can be represented as a deterministic
map
\[
g:\Sigma_{\mathrm{obs}}\to Z
\]
acting on the KPI readout \(\pi_K(\omega)\).

Deterministic post-processing cannot refine the partition induced by
\(\pi_K\):
if two states yield identical KPI vectors, then all deterministic functions of
those vectors are also identical.
Consequently, no increase in sampling frequency, no aggregation strategy, and no
deterministic statistical enrichment can restore identifiability once it is lost
at the level of KPI readout.

This is a standard result in identification theory: information that is not
present in the observation cannot be reconstructed by downstream processing
alone \cite{ljung1999,walter1982}.

\subsection{Interpretation of the Results}

The results established in this preprint are not claims about the usefulness,
quality, or desirability of KPIs.
They are strictly formal statements about limits of identifiability.

Within OC and \texttt{ETS.K\_EA.v1.1}, KPI-based observation modeled as a finite
readout map cannot be assumed to uniquely identify the underlying enterprise
state without introducing additional identifying structure or constraints that
lie outside the OC/ETS framework.

Read in this way, the results function as a clarification of scope: they specify
what KPI-based analysis can and cannot be expected to recover at the level of
formal enterprise continua.
